\thispagestyle{empty}\begingroup\fontsize{9}{11}\selectfont\raggedright\noindent{\Large\bfseries Übersetzungsquellen und Terminologie\par}\medskip\noindent Der lateinische Text stammt aus Carl Friedrich Gauss, \textit{Werke}, Band IV (Göttingen, 1873), herausgegeben von der Königlichen Gesellschaft der Wissenschaften zu Göttingen. Die originalsprachliche editierbare Transkription und das historische Porträt stehen in der \href{https://github.com/KokunoYumeto/gauss-collected-works/releases/tag/gauss-original-language-2026-10-24}{originalsprachlichen Ausgabe}. Die moderne Quellkodierung und der Satz sind OpenAI GPT-6 Astra, Ultra, zugeschrieben; diese deutsche Übersetzung und ihr Satz sind KI-produziert.\par\medskip\noindent Verwandter englischer Sprachgebrauch für mathematische Argumente wurde in James Caddall Moreheads und Adam Miller Hiltebeitels \textit{General Investigations of Curved Surfaces of 1827 and 1825} (1902), \href{https://www.gutenberg.org/ebooks/36856}{Project Gutenberg 36856, TeX-Ausgabe}, herangezogen. Verwandter deutscher Sprachgebrauch für Größen und mathematische Argumente wurde in H. Masers Übersetzung und Ausgabe \textit{Carl Friedrich Gauss’ Untersuchungen über höhere Arithmetik} (1889), \href{https://archive.org/details/carlfriedrichga00gausgoog}{historische digitalisierte Ausgabe}, herangezogen. Diese Werke sind keine Parallelübersetzungen der Beobachtungstheorie. Die deutsche Anzeige von Gauss vom 26. Februar 1821 im selben historischen Band wurde für verwandte Fehler- und Wahrscheinlichkeitsterminologie gelesen. Menschliche laufende Prosa und Formeln werden aus diesen Referenzen nicht übernommen. Historische Publikationstitel bleiben im Original. Genaue Versionen und URLs stehen in \texttt{CITATIONS.bib}.\par\endgroup\clearpage
